\section{Gauge fields, oriented stars, and the test block}
\label{sec:setup}

Let $K$ be a finite oriented cubical complex.  A lattice gauge field assigns
$U_e\in\SU$ to every oriented edge, with $U_{\bar e}=U_e^{-1}$.  A gauge
transformation $g:K^0\to\SU$ acts by
\[
  U_{xy}\longmapsto g_xU_{xy}g_y^{-1}.
\]
Fix an edge $\ell=(x,y)$.  Reorient each incident plaquette so that its
boundary starts with $\ell$.  Its holonomy then has the form
\[
  W_i=U_\ell S_i,
\]
where $S_i$ is the parallel transport along the complementary three-edge path
from $y$ back to $x$.  Therefore
\[
 S_i\longmapsto g_yS_i g_x^{-1},\qquad
 U_\ell^{-1}\longmapsto g_yU_\ell^{-1}g_x^{-1}.
\]
The $S_i$ do not transform by conjugation, but all transform by the same left
and right factors.  They occupy a common transport bitorsor.  Any expression
built only from their mutual bi-invariant distances is gauge invariant.

We use the bi-invariant metric for which $\SU$ is the unit round three-sphere.
Thus the distance from the identity is the quaternion class angle in
$[0,\pi]$.  Rescaling the metric rescales every quadratic form in the paper
and leaves all dimensionless spectral ratios unchanged.

For the exact calculations, $K=[0,2]^4\cap\mathbb Z^4$ with open boundary.
It has
\[
  |K^0|=81,\qquad |K^1|=216,\qquad |K^2|=216.
\]
Let
\[
 d_1:C^1(K;\R)\longrightarrow C^2(K;\R)
\]
be the signed link-to-plaquette coboundary.  It has rank $136$.  There are
eight interior links, and each belongs to six plaquettes.  The same scalar
matrix acts independently on the three color components in the linearized
$\mathfrak{su}(2)\simeq\R^3$ theory.

For an interior link $\ell$, let $S_\ell$ select its six incident plaquette
rows and let
\[
 O_\ell=\operatorname{diag}\bigl(d_{1,p\ell}:p\supset\ell\bigr).
\]
Multiplication by $O_\ell$ reorients all six rows so that the coefficient of
the distinguished link is $+1$.  Finally set
\[
 C_6=I_6-\frac16\mathbf1\mathbf1^\mathsf T.
\]
The order $C_6O_\ell S_\ell$ is not cosmetic: without $O_\ell$, centering
signed plaquette coordinates does not linearize the staple dispersion.
